% =====================================================================
% grain_initialization_keff.tex -- how the grains are initialized for the
% large-domain runs that feed the effective-thermal-conductivity
% upscaling (studies/keff_sintering/).
%
%     pdflatex grain_initialization_keff.tex
%
% Sources of truth, in the order the pipeline runs them:
%   preprocess/generate_packing_gravity.py   deposition, fill, gates
%   preprocess/packing_lib.py                metrics, relax, grains.dat
%   preprocess/comp_eps.py                   eps and the mesh
%   preprocess/generate_study_opts.py        the .opts matrix
%   src/initial_conditions.c                 phi, T, rho_v at t = 0
% If any of those change, this changes with them.
% =====================================================================
\documentclass[11pt]{article}
\usepackage[margin=1in]{geometry}
\usepackage{amsmath, amssymb}
\usepackage{booktabs}
\usepackage{textcomp}
\usepackage[T1]{fontenc}

\newcommand{\code}[1]{\texttt{\small #1}}

% long \code{} tokens are unbreakable; give TeX room rather than overfull boxes
\setlength{\emergencystretch}{2em}

\title{Grain initialization for the $k_{\text{eff}}$ upscaling runs}
\author{enceladus\_DSM}
\date{2026-09-25}

\begin{document}
\maketitle

\noindent
The large-domain runs behind \code{studies/keff\_sintering/} start from a
\emph{constructed} granular microstructure, not from an equilibrated field
and not from a mechanical settling simulation. Initialization happens in
two clearly separated stages:

\begin{enumerate}
  \item \textbf{Offline (Python).} A sharp geometry --- a list of
    non-overlapping circles $(x_k, y_k, R_k)$ --- is generated,
    \emph{graded}, and committed to the repository as
    \code{inputs/packings/.../grains.dat}. No phase field, no $\epsilon$,
    no mesh is involved; this stage is temperature-independent and is
    reused at every run temperature.
  \item \textbf{Online (solver).} At $t=0$ the solver reads that list and
    paints a diffuse phase field over it, together with the temperature
    and vapour-density fields. This stage is where $\epsilon$ and the mesh
    enter.
\end{enumerate}

\noindent
Keeping the two apart is the design decision the rest of this note
elaborates: the sharp geometry is fixed and auditable, and $\epsilon$ only
sets the thickness of the diffuse skin painted around it.

\section{What the packings are meant to be}

Every $k_{\text{eff}}$ packing targets the same microstructure and differs
only in box size and random seed:

\begin{center}
\begin{tabular}{lll}
\toprule
porosity $\phi$            & $0.325$                    & gated to $\pm 0.01$ \\
mean grain radius $R_{ave}$ & $50\ \mu$m                & log-normal, $\sigma_{\ln} = 0.5$ \\
radius clip                & $[0,\,2R_{ave}]$           & \code{--radius-clip-frac 1.0} \\
periodicity                & both axes (\code{xy})      & the only mode with no boundary layer \\
dimension                  & 2                          & \\
\bottomrule
\end{tabular}
\end{center}

\noindent
Two box sizes are in use, both set by the RVE ratio $L/R_{ave}$ rather than
by a length chosen directly:
%
\begin{align*}
  \textbf{pilot:} && L/R_{ave} &= 40 && L = 2.0\ \text{mm}  && \text{\code{inputs/packings/pilot\_LR40/}} \\
  \textbf{REV:}   && L/R_{ave} &= 64 && L = 3.2\ \text{mm}  && \text{\code{inputs/packings/rev\_LR64/}}
\end{align*}
%
with four seeds each, so seed scatter is measured rather than assumed.
Periodicity in \emph{both} axes is required: under \code{-periodic 0} the
zero-flux wall imposes a $90^\circ$ contact angle on $\phi$, so a grain
clipped by a wall carries the wrong surface curvature and coarsens as
though it were smaller than it is.

\section{Stage 1 (offline): building the sharp packing}

\code{preprocess/generate\_packing\_gravity.py}. Five steps, in order.

\paragraph{1.1 Geometric drop-and-roll deposition.}
Grains are placed one at a time into a strip of width $L$. A radius is
drawn from the log-normal, an $x$ is drawn uniformly, and the grain is
lowered until it first touches the floor or a settled grain; it then
\emph{rolls downhill} around that contact until it finds a second support,
exhausts a rolling budget, reaches the floor, or rolls past the support's
equator (in which case it falls off and is re-dropped, up to four hops).

There is no DEM, no contact law and no time integration --- this must not
be reported as a settling simulation. What it buys is the one structural
property that matters here: every deposited grain rests on something below
it, so the solid matrix is connected \emph{by construction} rather than by
luck, and no placement that overlaps a settled grain is accepted.

\paragraph{1.2 Porosity by bisecting a physical knob.}
The rolling budget \code{--roll-tol} is the porosity control:
$0$ means lock on first touch (ballistic, $\phi \approx 0.45$--$0.50$) and
$\pi/2$ means roll to a stable two-point seat every time (dense,
$\phi \approx 0.16$--$0.20$). Porosity falls monotonically with the budget,
so it is bisected \emph{per seed} onto the target. Radii are never rescaled
to hit a porosity: rescaling would break contact tangency, which is the
whole point (see \emph{Why additive} in \S\ref{sec:online}). The realized budgets for the eight
packings in use run $0.94$--$1.26$ rad.

Deposition always wraps in $x$ regardless of the requested output
periodicity: the bed being modelled is wide, and a strip with free sides
opens spurious voids at the edges and drives the interior too dense.

\paragraph{1.3 The domain is a window cut from a larger bed.}
Deposition runs in a strip taller than the domain by
$3 R_{ave}$ above and below, and the output window is cropped from its
interior. Both discarded layers are artifacts: the floor contact layer at
the bottom and the loose, half-filled free surface at the top. Deposition
stops on the \emph{minimum} of the bed surface profile, not the highest
grain, so the crop never reaches up into that loose layer.

\paragraph{1.4 Seam healing, then filling --- in that order.}
Identifying $y=0$ with $y=L_y$ leaves grains near the seam overlapping, so
the seam is relaxed apart first (\code{packing\_lib.relax} under the
minimum image, with $x$ wrapped). Maximum grain displacement is recorded as
\code{seam\_max\_shift\_m}; it is $28$--$81\ \mu$m, i.e.\ confined to grains
at the seam.

The deposited skeleton is deliberately bisected to land $0.03$ \emph{above}
the target porosity, and the remainder is removed by inserting grains into
the largest voids, biggest hole first. Attacking the largest inscribed
empty circle serves two goals with one pass: porosity falls to target and
pore heterogeneity falls with it. Each insert is capped by the largest
circle that fits \emph{and} by the remaining area deficit, and the last one
is shrunk to land exactly on target rather than overshoot.

Filler grains are not required to rest on anything --- some float. That is
intentional: a 2D section through a 3D packing shows unsupported grains all
the time, and a grain connected out of plane is mechanically real even
where the slice cannot show it. The matrix is built connected \emph{first}
so the floaters garnish a percolating skeleton instead of constituting it,
and \code{metadata.json} records the split (\code{n\_matrix\_grains} vs
\code{n\_filler\_grains}: $4$--$6\%$ filler across the eight packings).

\paragraph{1.5 Acceptance is a gate, not a ranking.}
For homogenization the domain has to be a usable REV, so a packing that is
well connected on one side and sparse on the other is worse than useless:
it yields a number that looks fine and means nothing. The generator
therefore \emph{refuses} to emit a packing and re-seeds (up to 24 times)
unless all of
%
\begin{itemize}\itemsep2pt
  \item the largest inscribed pore circle is below its gate,
  \item the local density CV on a coarse grid is below $0.3475$,
  \item the half-domain solid-fraction difference is below $22\%$ in either
        axis,
  \item the solid percolates in \emph{both} $x$ and $y$,
  \item achieved porosity is within $0.01$ of target
\end{itemize}
%
hold. Three of the four pilot packings needed 2--3 attempts. Packings are
never \emph{ranked} on pore connectivity: maximising the fraction of void
in one cluster selects for clumped grains beside one large hole, which is
the opposite of representative.

The void gate is applied per domain size, not in mean radii: the largest of
$N$ voids grows with $N$, so a flat mean-radii threshold calibrated on the
$L/R_{ave}=40$ box silently rejects the larger and relatively \emph{more}
uniform $L/R_{ave}=64$ box.\footnote{Visible in the metadata: REV seed~1
records \code{gate\_max\_void\_ratio} $=2.14 = 1.3375\times 64/40$, i.e.\
the default rescaled to the larger box, while seeds 2--4 record the unused
default beside void ratios above it, i.e.\ they were gated on
\code{--max-void-per-L}. Both are the same size-independent criterion, and
all four REV packings sit at $2.1$--$2.5\%$ of $L$ against the
$3.3$--$3.4\%$ the accepted pilot packings reach.}

\paragraph{1.6 What is written out.}
\code{grains.dat} (an \code{Lx Ly} header then \code{x y r} rows, all in
metres), \code{metadata.json} (every graded metric, the gates, the seed and
attempt, and a SHA-256 of \code{grains.dat}) and \code{preview.png}.

One detail in \code{grains.dat} matters more than its size suggests:
\textbf{periodic edge images are written explicitly}. Any grain reaching
past a periodic face is emitted a second time at its image, so centres span
$[-R,\,L+R]$ and a grain need only be \emph{partly} inside the domain.
Dropping the edge-straddling grains would sever exactly the conduction
paths that leave the domain, and connectivity across the boundary is what
$k_{\text{eff}}$ is most sensitive to --- it would bias $k_{\text{eff}}$
low. The images also make the microstructure genuinely periodic across
those faces, which is what the corrector problem's periodic boundary
conditions assume. The pilot files carry $25$--$30$ images on top of
$308$--$336$ distinct grains.

\section{Stage 2 (online): the fields at $t = 0$}
\label{sec:online}

\code{src/initial\_conditions.c}, reached via \code{-ic\_type
multi\_grains\_file} and \code{-grains\_file <path>}. The grain list is
parsed on rank 0 and broadcast, so every rank sees byte-identical geometry;
a ragged or zero-radius file is an error rather than a silent
misinterpretation.

\paragraph{2.1 The ice phase field: additive, clamped $\tanh$ profiles.}
$\phi$ is the clamped \emph{sum} of the grains' individual diffuse fields
(\code{-ic\_grain\_union 0}, the Molaro convention):
%
\begin{equation}
  \phi(\mathbf{x}) \;=\;
  \operatorname{clamp}\!\left(
    \sum_{k}
      \left[\tfrac12 - \tfrac12 \tanh\!\left(\frac{r_k - R_k}{2\epsilon}\right)\right],
    \;0,\;1\right),
  \qquad r_k = \lVert \mathbf{x} - \mathbf{c}_k \rVert .
  \label{eq:phi}
\end{equation}
%
The $\tanh$ coefficient is $1/(2\epsilon)$, i.e.\ the model's own
equilibrium logistic profile $\phi = 1/(1+e^{-x/\epsilon})$. Initializing
at the equilibrium width removes the early width-relaxation transient that
a steeper initial profile would spend its first $\sim 60$ steps undoing.
Note that $\epsilon$ is a \emph{decay length}: the visible band is
$\approx 6\epsilon$ ($5\%$--$95\%$) to $\approx 9.2\epsilon$
($1\%$--$99\%$), not $\epsilon$.

\paragraph{2.2 Why additive, and why contact is $\epsilon$-independent.}
Each grain's own term in \eqref{eq:phi} crosses $0.5$ at exactly
$r_k = R_k$ --- the sharp grain surface --- \emph{for any} $\epsilon$. Since
drop-and-roll seats every deposited grain in exact tangency on its
supports, neighbours touch precisely where $\phi_1 = \phi_2 = 0.5$.
Three consequences:
%
\begin{itemize}\itemsep2pt
  \item contact has an $\epsilon$-independent meaning;
  \item the radii never depend on $\epsilon$, so the same
        \code{grains.dat} is valid at every temperature and resolution;
  \item no neck calibration step is needed --- the sharp geometry is fixed
        and only the skin thickness changes.
\end{itemize}
%
The alternative form, $\tanh$ of the signed distance to the grains' union
(\code{-ic\_grain\_union 1}), exists for \emph{overlapping} grains. These
packings deliberately have none, so it is not used here.

\paragraph{2.3 Temperature and vapour.}
Every IC shares one closure: temperature is the linear background field
$T = T_0 + \nabla T_0 \cdot (\mathbf{x} - \mathbf{L}/2)$ (the
$k_{\text{eff}}$ campaign sets no gradient, so $T \equiv T_0$), and vapour
density is saturated inside ice and $h_0$-scaled in air, blended by the
ice fraction:
%
\begin{equation}
  \rho_v = \rho_{vs}(T)\,\big[\,h_0\,a + (1-a)\,\big],
  \qquad a = \max(0,\,1-\phi),
\end{equation}
%
with $h_0 = 1$ (saturated) for these runs.

\paragraph{2.4 Where $\phi$ is evaluated.}
Node coordinates come from the Greville abscissae of the B-spline basis,
with the periodic-axis node spacing using $N$ intervals rather than $N-1$
(a periodic node at $x=L$ coincides with the one at $x=0$). Getting that
denominator wrong stretched every IC by $(N+1)/N$ and left a double-width
gap at the seam --- measured as a $2/N$ disc-area excess in
\code{studies/keff\_sharp\_limit/}.

\paragraph{2.5 $\epsilon$ and the mesh.}
$\epsilon$ is not chosen by eye. \code{preprocess/comp\_eps.py} takes the
minimum of the three Kaempfer \& Plapp (2009) sharp-interface bounds (with
safety $s=\tfrac12$) and the geometric bound $0.05\,R_{ave}$, with the
normal velocity in the kinetic bound derived self-consistently from a
feature size $R_{feat} = 2\ \mu$m so that the bound cannot collapse at cold
temperatures. At $T=-20\,^\circ$C this gives $\epsilon = 1\ \mu$m for these
packings, with the heat bound binding. The mesh follows:
%
\begin{equation}
  h = \frac{\epsilon}{\sqrt{2}}, \qquad
  N_x = \left\lceil \frac{\sqrt{2}\,L_x}{\epsilon} \right\rceil,
\end{equation}
%
i.e.\ $2829^2$ for the pilot and $4526^2$ for the REV box, i.e.\ $24$M and
$61$M degrees of freedom at $3$ fields per node. The geometry \code{.opts}
also carries \code{-eps\_valid\_temp}, and the solver aborts if the run
temperature disagrees with it by more than $1\,^\circ$C --- $\epsilon$ and
the mesh are temperature-specific even though the packing is not.

\section{What was actually built}

Realized values from \code{metadata.json}, all at $\phi_{target}=0.325$,
$R_{ave}=50\ \mu$m, \code{periodic xy}:

\begin{center}
\small
\begin{tabular}{llrrrrrrr}
\toprule
set & seed & grains & filler & $\phi$ & $Z$ & $Z_{band}$ & roll [rad] & rows \\
\midrule
pilot, $L=2.0$\,mm   & 1 & 309 & 14 & 0.3251 & 1.96 & 3.40 & 0.957 & 334 \\
                     & 2 & 308 & 15 & 0.3251 & 1.99 & 3.49 & 0.982 & 337 \\
                     & 3 & 336 & 21 & 0.3249 & 1.99 & 3.41 & 1.006 & 366 \\
                     & 4 & 310 & 20 & 0.3252 & 1.92 & 3.36 & 1.080 & 336 \\
\midrule
REV, $L=3.2$\,mm     & 1 & 797 & 27 & 0.3248 & 1.98 & 3.57 & 1.104 & 845 \\
                     & 2 & 805 & 32 & 0.3254 & 2.03 & 3.53 & 1.115 & 844 \\
                     & 3 & 724 & 37 & 0.3249 & 1.93 & 3.52 & 1.264 & 758 \\
                     & 4 & 769 & 42 & 0.3254 & 1.95 & 3.33 & 0.945 & 815 \\
\bottomrule
\end{tabular}
\end{center}

\noindent
\code{rows} is the line count in \code{grains.dat}, i.e.\ distinct grains
plus periodic edge images. Two coordination numbers are reported because
they answer different questions: $Z$ counts near-exact contacts
($\approx 2$, barely rigid), while $Z_{band}$ counts every pair closer than
the diffuse band $9.2\epsilon$ ($\approx 3.4$, nearly isostatic).
\textbf{$Z_{band}$ is the number that describes the simulation}, because
the solver joins with solid every pair closer than the band. The solid
percolates in both axes in all eight packings (largest cluster
$75$--$96\%$ of the solid).

\section{Two consequences to keep in view}

\paragraph{The diffuse band closes narrow throats at $t=0$.}
Because \eqref{eq:phi} is additive, two grains that nearly touch each
contribute $\approx 0.5$ at the midpoint of the gap and \emph{sum} to
$\approx 1$. At the centre of a gap $g$ the sum is
$1 - \tanh\!\big(g/4\epsilon\big)$, so $\phi$ only falls below $0.01$ once
%
\begin{equation}
  g > 4\epsilon \operatorname{artanh}(0.99) \approx 10.6\,\epsilon .
\end{equation}
%
At $\epsilon = 1\ \mu$m that is $10.6\ \mu$m, against a median pore throat
of $19.4\ \mu$m (pilot) / $15.6\ \mu$m (REV) and a $25$th percentile of
$\approx 1.5\ \mu$m. So roughly a quarter of the throats hold no pore at
all at $t=0$, and a channel narrower than the band conducts like solid
however open the drawing looks. This is recorded in the header of every
geometry \code{.opts} and is why the deliverable of the campaign is
\emph{comparative} --- how much sintering changes $k_{\text{eff}}$ --- and
not absolute $k_{\text{eff}}$ in W\,m$^{-1}$\,K$^{-1}$. The same effect
puts the usable (band-eroded) pore space in hundreds of disconnected
pockets at $t=0$, which bears on vapour transport rather than on
conduction.

\paragraph{$t=0$ is not the baseline; $t=1$\,day is.}
Equation \eqref{eq:phi} is an analytic clamped sum, not an equilibrated
phase field, and the grains start in exact \emph{point} tangency with no
neck. The first hours of every run are therefore the field relaxing rather
than the microstructure sintering, and it is precisely at $t=0$ that a band
bridging a point contact does the most damage to $k_{\text{eff}}$: on pilot
seed~1, $k_{\text{eff}}(0)$ reads $0.6312$ under the arithmetic
conductivity law against $0.3695$ under the tensor law, i.e.\ the
arithmetic law was $+71\%$ high on tangent grains. All reported changes are
measured against the $t = 1$~day state.

\section{Reproducing it}

The generation commands are not committed; the following are reconstructed
from \code{metadata.json} and produce the committed packings (the generator
is deterministic in the seed). The flag
\code{--no-periodic-subdir} is what puts the output at the committed path
rather than under a \code{periodic\_xy/} level.

\begin{quote}\small\ttfamily
python3 preprocess/generate\_packing\_gravity.py \textbackslash \\
\hspace*{2em} --Lx 2.0e-3 --porosity 0.325 --mean-r 50e-6 --sigma-ln 0.5 \textbackslash \\
\hspace*{2em} --periodic xy --band-per-mean-r 0.184 --seed 1 \textbackslash \\
\hspace*{2em} --no-periodic-subdir \textbackslash \\
\hspace*{2em} --out inputs/packings/pilot\_LR40/pilot\_phi0.325\_Rave50um\_LR40\_seed1
\end{quote}

\noindent
For the REV set, \code{--Lx 3.2e-3} and the size-independent void gate
\code{--max-void-per-L 0.0335}. The \code{.opts} matrix on top of a packing
comes from \code{preprocess/generate\_study\_opts.py} with
\code{--R\_feat 2e-6}, and runs go through the run scripts
(\code{scripts/HPC/submit\_enceladus.sh <geometry> <experiment> [tag]}),
never the bare executable, so the staged inputs and figures exist
afterwards.

\paragraph{Seeing the initial condition the solver actually uses.}
The generator's \code{preview.png} draws sharp discs, which is the geometry
as specified rather than the field the run starts from. To render
\eqref{eq:phi} itself, including a zoom far enough in to make the
$9.2\ \mu$m band visible at the $2$\,mm domain scale (it is $0.46\%$ of the
frame):

\begin{quote}\small\ttfamily
venv\_enceladus/bin/python studies/keff\_sintering/render\_pilot\_ic.py
\end{quote}

\end{document}
